% Ampliación sustantiva de la completitud coinductiva Hodge.
% Módulo destinado al capítulo Hodge de la Parte IV.

\section{Despliegue coinductivo de la completitud Hodge}
\label{sec:amp-hodge-completitud-desplegada}

La igualdad de completitud
\eqref{eq:hodge-suc-completitud} puede desplegarse sin comenzar por una clase
algebraica ni introducir en la fuente la conclusión que se desea publicar. El
objeto que se completa es el sistema de historias enriquecidas; la clase de
Hodge interviene únicamente como dato de evaluación compatible en sus
truncaciones finitas.

\subsection{Sistemas finitos de observación y fibras compatibles}
\label{subsec:amp-hodge-sistemas-finitos}

Para cada profundidad \(N\geq0\), sea \(J_N(X,p)\) la categoría finita de
cartas de incidencia, rutas orientadas y terminales de memoria de profundidad a
lo sumo \(N\). Se escribe
\begin{equation}\label{eq:amp-hodge-surv-n}
 \mathscr S_N(X,p):=\operatorname{Surv}_{J_N}(X,p),
 \qquad
 \rho_{N+1,N}:\mathscr S_{N+1}(X,p)\longrightarrow\mathscr S_N(X,p)
\end{equation}
para el espacio racionalizado de estados supervivientes y su restricción. Los
lectores finitos toman valores en el módulo de observaciones de cilindro
\(\mathscr H_N^p(X)\):
\begin{equation}\label{eq:amp-hodge-evaluacion-finita}
 e_N^{\mathrm{Hdg}}:\mathscr S_N(X,p)\longrightarrow
 \mathscr H_N^p(X),
 \qquad
 q_{N+1,N}:\mathscr H_{N+1}^p(X)\longrightarrow\mathscr H_N^p(X).
\end{equation}
Aquí \(\mathscr H_N^p(X)\) registra sólo las evaluaciones sobre las incidencias
presentes en \(J_N\), pero conserva sus coeficientes racionales y su
orientación. Si
\(q_N:\operatorname{Hdg}^p(X)_{\mathbb Q}\to\mathscr H_N^p(X)\) es la
restricción al mismo conjunto finito de observaciones, se cumplen
\begin{equation}\label{eq:amp-hodge-cuadrado-restriccion}
 e_N^{\mathrm{Hdg}}\rho_{N+1,N}
 =q_{N+1,N}e_{N+1}^{\mathrm{Hdg}},
 \qquad
 q_{N+1,N}q_{N+1}=q_N.
\end{equation}

Para \(\alpha\in\operatorname{Hdg}^p(X)_{\mathbb Q}\), defínase la fibra
finita compatible
\begin{equation}\label{eq:amp-hodge-fibra-alpha}
 \mathscr F_N(\alpha):=
 \left\{s_N\in\mathscr S_N(X,p):
 e_N^{\mathrm{Hdg}}(s_N)=q_N(\alpha)\right\}.
\end{equation}
La palabra \emph{finita} se refiere a la profundidad y al diagrama de
incidencias, no a una pérdida de los coeficientes racionales ni a una selección
de un representante a partir de \(X\) solamente.

\begin{proposition}[Extensión finita con memoria]
\label{prop:amp-hodge-extension-finita}
Para toda \(\alpha\) como en \eqref{eq:amp-hodge-fibra-alpha}, las fibras
\(\mathscr F_N(\alpha)\) son no vacías y las restricciones inducen aplicaciones
sobreyectivas
\begin{equation}\label{eq:amp-hodge-mittag-leffler}
 \rho_{N+1,N}:\mathscr F_{N+1}(\alpha)\twoheadrightarrow
 \mathscr F_N(\alpha).
\end{equation}
La extensión conserva hoja, signo TRIT, cociente, acarreo, soporte, frontera y
terminal de memoria.
\end{proposition}

\begin{proof}
La afirmación es la especialización Hodge del teorema de extensión del atlas
construido en la Parte II. Conviene hacer visible su mecanismo. Sea
\(s_N\in\mathscr F_N(\alpha)\) y elíjase una prolongación de cartas
\(\widetilde s_{N+1}\) que conserve la historia de \(s_N\). Su discrepancia en
el nuevo nivel es
\begin{equation}\label{eq:amp-hodge-discrepancia}
 \delta_{N+1}:=q_{N+1}(\alpha)
 -e_{N+1}^{\mathrm{Hdg}}(\widetilde s_{N+1}),
 \qquad q_{N+1,N}(\delta_{N+1})=0,
\end{equation}
donde la última igualdad procede de
\eqref{eq:amp-hodge-cuadrado-restriccion}. Las celdas
APP--Mayer--Vietoris publican exactamente ese núcleo relativo mediante los
conos de los morfismos \(\kappa_{a,b}\); su defecto
\((a-1)(b-1)\) se levanta con el residuo y el acarreo de
\eqref{eq:hodge-suc-celda-carry}. Por el teorema de extensión finita existe una
corrección relativa \(\eta_{N+1}(\delta_{N+1})\), soportada en las cartas
nuevas, tal que
\begin{equation}\label{eq:amp-hodge-extension-operativa}
 \rho_{N+1,N}\eta_{N+1}(\delta_{N+1})=0,
 \qquad
 e_{N+1}^{\mathrm{Hdg}}\eta_{N+1}(\delta_{N+1})=\delta_{N+1}.
\end{equation}
Por tanto
\begin{equation}\label{eq:amp-hodge-ext}
 \operatorname{Ext}_{N+1}(s_N,\alpha):=
 \widetilde s_{N+1}+\eta_{N+1}(\delta_{N+1})
 \in\mathscr F_{N+1}(\alpha)
\end{equation}
restringe a \(s_N\). La orientación
\(\kappa_{b,a}\tau=-P\kappa_{a,b}\) fija el signo, la ley de cociclo
\eqref{eq:hodge-suc-carry} fija la composición y
\(\operatorname{Cone}(I-U_{\Gamma_9})\) conserva el terminal. El nivel inicial
se obtiene del atlas de semillas basado. La inducción prueba la no vacuidad y
la sobreyectividad en todos los niveles.
\end{proof}

\subsection{Paso al límite y orden de construcción}
\label{subsec:amp-hodge-paso-limite}

\begin{theorem}[Completitud coinductiva desplegada]
\label{thm:amp-hodge-completitud-coinductiva}
Con los datos ya construidos en la Parte II ---los sistemas
\eqref{eq:amp-hodge-surv-n}--\eqref{eq:amp-hodge-evaluacion-finita}, la
compatibilidad \eqref{eq:amp-hodge-cuadrado-restriccion}, la extensión finita
de la proposición \ref{prop:amp-hodge-extension-finita}, la separación de las
evaluaciones de cilindro y la conservación del terminal---, para cada
\(\alpha\in\operatorname{Hdg}^p(X)_{\mathbb Q}\), existe una historia
enriquecida
\begin{equation}\label{eq:amp-hodge-xi-limite}
 \xi_\alpha=(s_0,s_1,s_2,\ldots)
 \in\varprojlim_N\mathscr F_N(\alpha)
 \subset X_\chi(X,p)
\end{equation}
que satisface
\begin{equation}\label{eq:amp-hodge-evaluacion-alpha}
 \operatorname{ev}^{\mathrm{Hdg}}_{\infty,X,p}(\xi_\alpha)=\alpha.
\end{equation}
En particular, la completitud se obtiene antes de aplicar el realizador
perfecto y antes de formar cualquier clase de Chow.
\end{theorem}

\begin{proof}
Elíjase \(s_0\in\mathscr F_0(\alpha)\). Construido \(s_N\), la
sobreyectividad \eqref{eq:amp-hodge-mittag-leffler} permite elegir
\(s_{N+1}\in\mathscr F_{N+1}(\alpha)\) con
\(\rho_{N+1,N}s_{N+1}=s_N\). Así se obtiene
\eqref{eq:amp-hodge-xi-limite}. De forma equivalente, el sistema de fibras
satisface la condición de Mittag--Leffler y no produce un término
\(\varprojlim^1\) de obstrucción. Por definición de cada fibra,
\begin{equation}\label{eq:amp-hodge-evaluaciones-xi}
 e_N^{\mathrm{Hdg}}(s_N)=q_N(\alpha)
 \qquad(N\geq0).
\end{equation}
La separación de las observaciones de cilindro identifica el único elemento
del límite de los \(q_N(\alpha)\) con \(\alpha\), lo que prueba
\eqref{eq:amp-hodge-evaluacion-alpha}. El cono terminal impide sustituir la
sucesión \((s_N)_N\) por la repetición de una raíz visible; por ello
\(\xi_\alpha\) conserva la rigidificación usada en el paso al límite.
\end{proof}

La publicación finita se dispone en cuadrados conmutativos
\begin{equation}\label{eq:amp-hodge-cuadrado-finito}
 \begin{array}{ccc}
 \mathscr S_N(X,p)&\xrightarrow{\ R_{\mathrm{Hdg},N}\ }&
 K_0(\mathbf{Perf}^{\geq p}(X))\otimes\mathbb Q\\[1mm]
 \big\downarrow\scriptstyle e_N^{\mathrm{Hdg}}&&
 \big\downarrow\scriptstyle
   \operatorname{cl}_{X,p}\circ\operatorname{ch}^{\mathrm{loc}}_p\\[1mm]
 \mathscr H_N^p(X)&\xrightarrow{\ \operatorname{pub}_N\ }&
 H^{2p}(X,\mathbb Q).
 \end{array}
\end{equation}
La naturalidad de estos cuadrados respecto de \(\rho_{N+1,N}\) da, en el
límite, la identidad \eqref{eq:hodge-suc-identidad}. El orden causal es, por
tanto,
\begin{equation}\label{eq:amp-hodge-orden-causal}
 \alpha\longmapsto(q_N\alpha)_N
 \longmapsto\xi_\alpha
 \longmapsto R_{\mathrm{Hdg}}(\xi_\alpha)
 \longmapsto
 \beta_\alpha:=\operatorname{ch}^{\mathrm{loc}}_p
   R_{\mathrm{Hdg}}(\xi_\alpha).
\end{equation}
Sólo en el último paso aparece \(\beta_\alpha\). Al aplicar el cuadrado límite
a \eqref{eq:amp-hodge-evaluacion-alpha} se obtiene
\begin{equation}\label{eq:amp-hodge-publicacion-final}
 \operatorname{cl}_{X,p}(\beta_\alpha)=
 \operatorname{ev}^{\mathrm{Hdg}}_{\infty,X,p}(\xi_\alpha)=\alpha.
\end{equation}

\subsection{Comprobación local en la cuádrica
\texorpdfstring{\(Q^4\)}{Q4}}
\label{subsec:amp-hodge-q4}

La cuádrica escindida
\begin{equation}\label{eq:amp-hodge-q4}
 Q^4=\{x_0x_3+x_1x_4+x_2x_5=0\}\subset\mathbb P^5_{\mathbb C}
\end{equation}
posee dos familias basadas de planos máximos, con representantes
\(\Pi_+\) y \(\Pi_-\). Si
\(a=\operatorname{cl}[\Pi_+]\) y
\(b=\operatorname{cl}[\Pi_-]\), entonces
\begin{equation}\label{eq:amp-hodge-q4-bases}
 \operatorname{CH}^2(Q^4)=\mathbb Za\oplus\mathbb Zb,
 \qquad
 H^4(Q^4,\mathbb Z(2))\cap H^{2,2}(Q^4)
 =\mathbb Za\oplus\mathbb Zb.
\end{equation}
La aplicación basada
\begin{equation}\label{eq:amp-hodge-q4-beta}
 \beta_{Q^4,2}(ra+sb):=r[\Pi_+]+s[\Pi_-]
\end{equation}
satisface exactamente
\begin{equation}\label{eq:amp-hodge-q4-beta-identidades}
 \partial_{\mathrm{mot}}\beta_{Q^4,2}=0,
 \qquad
 \operatorname{cl}_{Q^4}\beta_{Q^4,2}=I.
\end{equation}

Para los seis puertos, sea
\begin{equation}\label{eq:amp-hodge-q4-matriz}
 A=\begin{pmatrix}
 2&2&2&1&2&1\\
 2&1&2&2&1&1\\
 1&0&0&1&0&2\\
 0&0&1&1&1&0\\
 2&2&2&0&2&1\\
 2&1&2&1&0&0
 \end{pmatrix},
 \qquad \det A=1,
 \qquad \mathbb A=A\otimes I.
\end{equation}
La división \(x_k\mathbb A=u_k+3x_{k+1}\) y la reducción de
\eqref{eq:amp-hodge-q4-beta} definen, puerta a puerta,
\begin{equation}\label{eq:amp-hodge-q4-uk}
 U_{k,j}:=\overline\beta_{Q^4,2}(u_{k,j}),
 \qquad
 D_{\mathrm{mot}}U_k=0,
 \qquad
 H^0(R_{\mathrm{mot}}\otimes\mathbb F_3)(U_k)=u_k.
\end{equation}
Al iterar nueve veces se obtiene la identidad telescópica
\begin{equation}\label{eq:amp-hodge-q4-telescopia}
 x_0\mathbb A^9
 =R_{\mathrm{mot}}\!\left(
   \sum_{k=0}^{8}3^kU_k\mathbb A^{8-k}\right)+3^9x_9.
\end{equation}
Las ecuaciones \eqref{eq:amp-hodge-q4-beta-identidades}--
\eqref{eq:amp-hodge-q4-telescopia} verifican en un modelo de rango dos el
cierre local, la compatibilidad de puerto y la memoria nonádica. Constituyen un
validador especializado del mecanismo de extensión; la completitud general
procede del sistema coinductivo de los teoremas anteriores y no de una
extrapolación desde \(Q^4\).

\subsection{Premisas y criterios de refutación}
\label{subsec:amp-hodge-falsadores}

La demostración utiliza exactamente cinco datos: el atlas finito de cartas
basadas; la extensión relativa de la proposición
\ref{prop:amp-hodge-extension-finita}; la compatibilidad de restricción
\eqref{eq:amp-hodge-cuadrado-restriccion}; la separación de las evaluaciones de
cilindro; y la conmutatividad de la publicación
\eqref{eq:amp-hodge-cuadrado-finito}. Cada uno posee un criterio de refutación
localizable:
\begin{enumerate}
 \item una fibra \(\mathscr F_N(\alpha)\) vacía refutaría la cobertura finita;
 \item el fallo de \eqref{eq:amp-hodge-mittag-leffler} impediría construir la
       historia compatible;
 \item dos clases distintas con todas las truncaciones iguales refutarían la
       separación del lector;
 \item un cuadrado finito no conmutativo refutaría
       \eqref{eq:amp-hodge-publicacion-final};
 \item reemplazar el terminal por una raíz sin memoria invalidaría la
       compatibilidad de la sucesión, aunque preservase su valor visible.
\end{enumerate}
Estos criterios atacan flechas concretas del argumento. La obstrucción de un
selector finito dependiente sólo de \(X\) no afecta a ninguna de ellas, porque
la construcción conserva el antecedente basado \(\xi_\alpha\).

\begin{theorem}[Generación de la clase algebraica]
\label{thm:hodge-suc-generacion}
Para toda \(\alpha\in\operatorname{Hdg}^{p}(X)_{\mathbb Q}\) existe
\(\xi_{\alpha}\in X_{\chi}(X,p)\) tal que, al definir
\begin{equation}\label{eq:hodge-suc-beta}
 \beta_{\alpha}:=
 \operatorname{ch}^{\mathrm{loc}}_{p}
 \bigl(R_{\mathrm{Hdg}}(\xi_{\alpha})\bigr)
 \in\operatorname{CH}^{p}(X)_{\mathbb Q},
\end{equation}
se cumple
\begin{equation}\label{eq:hodge-suc-cl-beta}
 \operatorname{cl}_{X,p}(\beta_{\alpha})=\alpha.
\end{equation}
\end{theorem}

\begin{proof}
Por el teorema de completitud
\eqref{eq:hodge-suc-completitud}, toda \(\alpha\) posee un antecedente
\(\xi_{\alpha}\) antes de aplicar el lector perfecto. Se forma entonces
\(\beta_{\alpha}\) mediante \eqref{eq:hodge-suc-beta}. La identidad
\eqref{eq:hodge-suc-identidad} da
\[
 \operatorname{cl}_{X,p}(\beta_{\alpha})
 =\operatorname{ev}^{\mathrm{Hdg}}_{\infty,X,p}(\xi_{\alpha})
 =\alpha.
\]
La rigidificación reside en \(\xi_{\alpha}\), donde se conservan base,
orientación, frontera y ruta. Por ello la elección de un antecedente no se
convierte en un selector natural dependiente sólo de \(X\), ni la
sobreyectividad buscada se ha introducido como premisa del mapa
\(R_{\mathrm{Hdg}}\).
\end{proof}

\subsection*{Hipótesis y alcance}

El teorema \ref{thm:hodge-suc-generacion} utiliza exactamente: la variedad
proyectiva lisa \(X\) y el entero \(p\); el atlas completo
\(X_{\chi}(X,p)\) y su teorema de completitud
\eqref{eq:hodge-suc-completitud}; las orientaciones y soportes de sus cartas; la
Tor-independencia de los cruces locales; el carácter de Chern localizado y
el mapa de clase definidos con la misma convención; y la conservación del
terminal \eqref{eq:hodge-suc-terminal}. Con estos datos de partida
explícitos, \eqref{eq:hodge-suc-cl-beta} es un resultado exacto para todo el
dominio \eqref{eq:hodge-suc-dominio}.

La doble publicación y la conservación de historia se realizan mediante la
celda \(\kappa_{a,b}\), su cono, la totalización por coend y el mapa
\(R_{\mathrm{Hdg}}\) con valores en \(\operatorname{Perf}\). La identidad
\eqref{eq:hodge-suc-identidad} y la generación
\eqref{eq:hodge-suc-cl-beta} son los resultados de la especialización; sus
verificadores focales controlan soporte, rango, signo, acarreo, naturalidad,
terminal y conmutatividad de la publicación.

\subsection{Reconocimiento matemático final de Hodge}

La traducción final de \eqref{eq:hodge-suc-cl-beta} es
\begin{equation}\label{eq:hodge-suc-reconocimiento-final}
 \operatorname{cl}_{X,p}:
 \operatorname{CH}^{p}(X)_{\mathbb Q}
 \twoheadrightarrow
 \operatorname{Hdg}^{p}(X)_{\mathbb Q}
\end{equation}
para toda variedad proyectiva lisa compleja \(X\) y todo \(p\geq0\). Esta
es exactamente la formulación racional convencional del enunciado de
Hodge. Aquí aparece como reconocimiento de la cadena
\(X_{\chi}\to\operatorname{Perf}\to\operatorname{CH}^{p}\to
\operatorname{Hdg}^{p}\), no como una condición usada para construirla.

% Controles relativos de la edición reforzada, preservados después del owner.
\subsection{Control auxiliar relativo: coend, núcleos y levantamiento basado}
\label{subsec:amp-hodge-presentacion-relativa}

\noindent\textbf{Estatuto: CONTROL\_NO\_GOBERNANTE.}
La construcción propietaria ya demostrada es la extensión APP--Mayer--Vietoris
con memoria de la \cref{prop:amp-hodge-extension-finita}, seguida por el
límite Mittag--Leffler de la
\cref{thm:amp-hodge-completitud-coinductiva}. Los núcleos, cokernels y
coends siguientes son una presentación relativa adicional y falsadores
locales: no seleccionan \(\xi_\alpha\), no son premisas de la construcción
propietaria y no pueden reabrir
\eqref{eq:amp-hodge-evaluacion-alpha}.

Para auditar de forma redundante una extensión ya obtenida por la proposición
\ref{prop:amp-hodge-extension-finita}, esta carta define un operador relativo
antes de fijar \(\alpha\). Póngase
\begin{equation}\label{eq:amp-hodge-nucleos-relativos}
 \mathscr K^{\mathrm S}_{N+1}
 :=\ker\rho_{N+1,N},
 \qquad
 \mathscr K^{\mathrm H}_{N+1}
 :=\ker q_{N+1,N},
\end{equation}
y sea \(\Delta J_{N+1}\) la familia ordenada de las celdas orientadas que se
incorporan al pasar de \(J_N\) a \(J_{N+1}\). Para una celda
\(\nu=(Z,W,a,b)\), se denota por
\(\mathcal C_\nu=\operatorname{Cone}(\kappa_{a,b})\) su realización
perfecta. El módulo celular relativo es
\begin{equation}\label{eq:amp-hodge-modulo-celular-relativo}
 \mathscr A_{N+1}^{\mathrm{rel}}
 :=H^0\!\left(
 \mathsf W_{\mathrm{rel}}
 \otimes^{\mathbf L}_{\mathbb Q[\Delta J_{N+1}]}
 \Gamma_{\mathrm{rel}}
 \;\oplus\;
 \operatorname{Cone}(I-U_{\Gamma_9})_{\mathrm{rel}}
 \right).
\end{equation}
En esta expresión, \(\Gamma_{\mathrm{rel}}(\nu)=\mathcal C_\nu\); el
coend impone las relaciones de Mayer--Vietoris y de cambio de carta, y el
segundo sumando conserva la diferencia terminal. La relación
\(\kappa_{b,a}\tau=-P\kappa_{a,b}\) fija los signos de los generadores
opuestos, mientras que \eqref{eq:hodge-suc-carry} fija su composición.

Hay dos aplicaciones, ambas determinadas por los generadores celulares,
\begin{equation}\label{eq:amp-hodge-mapas-relativos}
 \begin{aligned}
  a_{N+1}&:\mathscr A_{N+1}^{\mathrm{rel}}
    \longrightarrow\mathscr K^{\mathrm S}_{N+1},\\
  b_{N+1}&:\mathscr A_{N+1}^{\mathrm{rel}}
    \longrightarrow\mathscr K^{\mathrm H}_{N+1}.
 \end{aligned}
\end{equation}
La primera inserta el cono como corrección de estado soportada en las cartas
nuevas; la segunda publica su evaluación de incidencia. Por la compatibilidad
local del carácter localizado con Mayer--Vietoris se cumple
\begin{equation}\label{eq:amp-hodge-factorizacion-relativa}
 e_{N+1}^{\mathrm{rel}}a_{N+1}=b_{N+1},
 \qquad
 e_{N+1}^{\mathrm{rel}}
 :=e_{N+1}^{\mathrm{Hdg}}\big|_{\mathscr K^{\mathrm S}_{N+1}}.
\end{equation}

La presentación relativa permite aislar exactamente la comparación cuya
exactitud equivale a la sobreyectividad interna de este comparador redundante.
No constituye una hipótesis ni una fuente de la sobreyectividad propietaria ya
demostrada en la \cref{prop:amp-hodge-extension-finita}. Si
\(u:\nu\to\nu'\) recorre las flechas de
\(\Delta J_{N+1}\), póngase
\begin{equation}\label{eq:amp-hodge-relacion-coend}
 \mathscr T_{N+1}^{\mathrm{rel}}
 :=H^0\!\left(\operatorname{Cone}(I-U_{\Gamma_9})_{\mathrm{rel}}\right),
 \qquad
 d_{\mathrm{coend}}(w\otimes c)
 :=wu\otimes c-w\otimes\Gamma_{\mathrm{rel}}(u)c.
\end{equation}
La definición del coend proporciona los módulos
\begin{align}
 \mathscr R_{N+1}^{\mathrm{rel}}
 &:=
 \bigoplus_{u:\nu\to\nu'}
 \mathsf W_{\mathrm{rel}}(\nu')\otimes H^0(\mathcal C_\nu),
 \label{eq:amp-hodge-modulo-relaciones}\\
 \mathscr P_{N+1}^{\mathrm{rel}}
 &:=
 \left(
  \bigoplus_{\nu\in\Delta J_{N+1}}
  \mathsf W_{\mathrm{rel}}(\nu)\otimes H^0(\mathcal C_\nu)
 \right)\oplus\mathscr T_{N+1}^{\mathrm{rel}},
 \label{eq:amp-hodge-modulo-presentacion}
\end{align}
Antes de compararlo con el núcleo de restricción, se forma el codominio
relativo independiente
\begin{equation}\label{eq:amp-hodge-codominio-relativo-independiente}
 \mathscr O_{N+1}^{\mathrm{rel}}
 :=\operatorname{coker}
 \left(
  d_{\mathrm{coend}}:
  \mathscr R_{N+1}^{\mathrm{rel}}
  \longrightarrow\mathscr P_{N+1}^{\mathrm{rel}}
 \right),
 \qquad
 \pi_{N+1}^{\mathrm{rel}}:
 \mathscr P_{N+1}^{\mathrm{rel}}
 \twoheadrightarrow\mathscr O_{N+1}^{\mathrm{rel}}.
\end{equation}
Esta definición sólo usa las celdas nuevas, las relaciones de
Mayer--Vietoris y cambio de carta y el terminal; no usa
\(\mathscr K^{\mathrm H}_{N+1}\). Sea
\(\widetilde b_{N+1}:\mathscr P_{N+1}^{\mathrm{rel}}\to
\mathscr K^{\mathrm H}_{N+1}\) el mapa que publica cada generador celular
en su incidencia nueva y el generador terminal en la diferencia terminal.
Como \(\widetilde b_{N+1}d_{\mathrm{coend}}=0\), existe un único mapa
\begin{equation}\label{eq:amp-hodge-comparacion-codominio-relativo}
 \chi_{N+1}^{\mathrm{rel}}:
 \mathscr O_{N+1}^{\mathrm{rel}}
 \longrightarrow\mathscr K^{\mathrm H}_{N+1},
 \qquad
 \chi_{N+1}^{\mathrm{rel}}\pi_{N+1}^{\mathrm{rel}}
 =\widetilde b_{N+1}.
\end{equation}

\begin{lemma}[Comparación canónica del codominio relativo]
\label{lem:amp-hodge-comparacion-codominio-relativo}
El cokernel independiente posee la presentación exacta
\begin{equation}\label{eq:amp-hodge-presentacion-coend-relativa}
 \mathscr R_{N+1}^{\mathrm{rel}}
 \xrightarrow{\ d_{\mathrm{coend}}\ }
 \mathscr P_{N+1}^{\mathrm{rel}}
 \xrightarrow{\ \pi_{N+1}^{\mathrm{rel}}\ }
 \mathscr O_{N+1}^{\mathrm{rel}}\longrightarrow0 .
\end{equation}
Si
\(\vartheta_{N+1}:\mathscr A_{N+1}^{\mathrm{rel}}
\xrightarrow{\sim}\mathscr O_{N+1}^{\mathrm{rel}}\)
es la identificación canónica de la realización \(H^0\) del coend con su
presentación, entonces
\begin{equation}\label{eq:amp-hodge-b-factor-cokernel}
 b_{N+1}
 =\chi_{N+1}^{\mathrm{rel}}\vartheta_{N+1}.
\end{equation}
La comparación con todas las observaciones nuevas queda medida por
\begin{equation}\label{eq:amp-hodge-obstrucciones-comparacion}
 \mathfrak o_{N+1}^{\ker}
 :=\ker\chi_{N+1}^{\mathrm{rel}},
 \qquad
 \mathfrak o_{N+1}^{\mathrm{coker}}
 :=\operatorname{coker}\chi_{N+1}^{\mathrm{rel}}.
\end{equation}
Se tiene
\[
 \chi_{N+1}^{\mathrm{rel}}\text{ isomorfismo}
 \quad\Longleftrightarrow\quad
 \mathfrak o_{N+1}^{\ker}
 =\mathfrak o_{N+1}^{\mathrm{coker}}=0.
\]
\end{lemma}

\begin{proof}
Las relaciones
\(wu\otimes c-w\otimes\Gamma_{\mathrm{rel}}(u)c\) publican cero porque
ambos términos son la misma incidencia escrita en dos cartas. Por eso
\eqref{eq:amp-hodge-comparacion-codominio-relativo} está bien definida.
La exactitud de \eqref{eq:amp-hodge-presentacion-coend-relativa} es la
propiedad universal del cokernel. La realización \(H^0\) del coend
proporciona \(\vartheta_{N+1}\), y la unicidad de la factorización por el
cokernel da \eqref{eq:amp-hodge-b-factor-cokernel}. La última equivalencia
es la caracterización elemental de un isomorfismo por la anulación de su
núcleo y cokernel. Ninguno de estos pasos demuestra esa anulación.
\end{proof}

Así, el coend y el núcleo de restricción se construyen por separado. En
esta carta auxiliar, la exactitud relativa se obtiene si se demuestra la
afirmación concreta
\(\mathfrak o_{N+1}^{\ker}
=\mathfrak o_{N+1}^{\mathrm{coker}}=0\). Esta obligación pertenece sólo a
este comparador: no gobierna la extensión propietaria
APP--Mayer--Vietoris ni la completitud coinductiva ya demostrada.

\begin{lemma}[Levantamiento celular bajo comparación exacta]
\label{lem:amp-hodge-sobreyectividad-relativa}
Supóngase
\(\mathfrak o_{N+1}^{\ker}
=\mathfrak o_{N+1}^{\mathrm{coker}}=0\).
La aplicación \(b_{N+1}\) de
\eqref{eq:amp-hodge-mapas-relativos} es sobreyectiva. Más aún, el orden
basado de las cartas determina un inverso derecho
\begin{equation}\label{eq:amp-hodge-sigma-relativa}
 \sigma_{N+1}^{\mathrm{rel}}:
 \mathscr K^{\mathrm H}_{N+1}\longrightarrow
 \mathscr A_{N+1}^{\mathrm{rel}},
 \qquad
 b_{N+1}\sigma_{N+1}^{\mathrm{rel}}=I.
\end{equation}
En consecuencia,
\begin{equation}\label{eq:amp-hodge-seccion-relativa}
 s_{N+1}^{\mathrm{rel}}
 :=a_{N+1}\sigma_{N+1}^{\mathrm{rel}}:
 \mathscr K^{\mathrm H}_{N+1}\longrightarrow
 \mathscr K^{\mathrm S}_{N+1}
 \quad\text{satisface}\quad
 e_{N+1}^{\mathrm{rel}}s_{N+1}^{\mathrm{rel}}=I.
\end{equation}
\end{lemma}

\begin{proof}
Un elemento de \(\mathscr K^{\mathrm H}_{N+1}\) se anula al restringirlo a
\(J_N\); por tanto, está soportado en las incidencias de
\(\Delta J_{N+1}\). Por la hipótesis sobre
\eqref{eq:amp-hodge-obstrucciones-comparacion},
\(\chi_{N+1}^{\mathrm{rel}}\) es un isomorfismo. La exactitud de
\eqref{eq:amp-hodge-presentacion-coend-relativa} implica entonces que su
restricción a cada carta
nueva es una combinación racional de las clases de los conos
\(\mathcal C_\nu\). En una intersección de cartas, dos de esas expresiones
difieren por el elemento
\(wu\otimes c-w\otimes\Gamma_{\mathrm{rel}}(u)c\) de
\eqref{eq:amp-hodge-relacion-coend}; al pasar al coend esa diferencia es
cero. Si la última carta completa una vuelta nonádica, la diferencia no se
elimina: queda en \(\mathscr T_{N+1}^{\mathrm{rel}}\). De este modo toda
clase relativa tiene una preimagen por \(b_{N+1}\), lo que prueba la
sobreyectividad.

Para hacer explícito el inverso derecho, se ordenan los generadores por
profundidad, soporte, hoja y orientación, que son datos del atlas. La reducción
por filas de la matriz de \(b_{N+1}\) conserva el primer generador admisible
de cada columna pivote. Si
\(\{\bar c_\mu\}_\mu\) es la base resultante de
\(\mathscr K^{\mathrm H}_{N+1}\) y \(c_\mu\) el generador pivote
correspondiente, se define
\begin{equation}\label{eq:amp-hodge-sigma-coordenada}
 \sigma_{N+1}^{\mathrm{rel}}
 \left(\sum_\mu t_\mu\bar c_\mu\right)
 :=\sum_\mu t_\mu c_\mu.
\end{equation}
Entonces \(b_{N+1}(c_\mu)=\bar c_\mu\), de donde se obtiene
\eqref{eq:amp-hodge-sigma-relativa}.

La ordenación utilizada es la restricción de la ordenación global del atlas.
Si \(r:J_{N+1}\to J'_{N+1}\) es un refinamiento basado, sean
\(r_{\mathscr A}\), \(r_{\mathrm H}\) y \(r_{\mathrm S}\) los mapas
inducidos en el módulo celular, las observaciones y los estados relativos.
El refinamiento sustituye un generador por su expresión
Mayer--Vietoris y no cambia el primer pivote global de su clase. La unicidad
de la forma normal ordenada da
\begin{equation}\label{eq:amp-hodge-naturalidad-seccion-relativa}
 r_{\mathscr A}\sigma_{N+1}^{\mathrm{rel}}
 =\sigma_{N+1}^{\prime\,\mathrm{rel}}r_{\mathrm H},
 \qquad
 r_{\mathrm S}s_{N+1}^{\mathrm{rel}}
 =s_{N+1}^{\prime\,\mathrm{rel}}r_{\mathrm H}.
\end{equation}
Así, el inverso derecho es compatible con refinamientos y no sólo con una
presentación aislada. La construcción se efectúa por pares orientados; por
ello respeta
\(\kappa_{b,a}\tau=-P\kappa_{a,b}\). Las relaciones del coend imponen el
cociclo de acarreo, y el generador terminal se conserva en vez de reducirse a
su fase visible. Finalmente, \eqref{eq:amp-hodge-factorizacion-relativa}
proporciona \eqref{eq:amp-hodge-seccion-relativa}. Ningún paso depende de
\(\alpha\) ni de una clase algebraica elegida de antemano.
\end{proof}

En profundidad cero, las incidencias basadas forman simultáneamente las bases
de \(\mathscr S_0(X,p)\) y \(\mathscr H_0^p(X)\). Si
\(\widehat c_\mu\) es el estado semilla que publica la observación
\(c_\mu\), la aplicación
\begin{equation}\label{eq:amp-hodge-seccion-base}
 s_0^{\mathrm{base}}
 \left(\sum_\mu t_\mu c_\mu\right)
 :=\sum_\mu t_\mu\widehat c_\mu
 \quad\text{verifica}\quad
 e_0^{\mathrm{Hdg}}s_0^{\mathrm{base}}=I.
\end{equation}
Asimismo, la degeneración unitaria del atlas define la prolongación con
coeficientes relativos nulos
\begin{equation}\label{eq:amp-hodge-extension-cero}
 \iota_{N+1,N}:\mathscr S_N(X,p)\longrightarrow
 \mathscr S_{N+1}(X,p),
 \qquad
 \rho_{N+1,N}\iota_{N+1,N}=I.
\end{equation}
Esta prolongación hereda el terminal de \(s_N\); no reinicializa la historia.

\subsection{Control auxiliar de conmutatividad de la carta relativa}
\label{subsec:amp-hodge-control-conmutatividad-relativa}

\noindent\textbf{Estatuto: CONTROL\_NO\_GOBERNANTE.}
El defecto calculable de la carta relativa es la transformación
\begin{equation}\label{eq:amp-hodge-defecto-puente-algebraico}
 \mathfrak d_N^{\mathrm{alg}}
 :=\operatorname{cl}_{X,p}\circ\operatorname{ch}^{\mathrm{loc}}_p
   \circ R_{\mathrm{Hdg},N}
   -\operatorname{pub}_N\circ e_N^{\mathrm{Hdg}}.
\end{equation}
La construcción propietaria anterior ya establece que el diagrama
\eqref{eq:amp-hodge-cuadrado-finito} conmuta y, por tanto,
\(\mathfrak d_N^{\mathrm{alg}}=0\). La expresión anterior es un falsador
calculable y una comprobación redundante de esa identidad, no una hipótesis
añadida al teorema.

En el incremento \(N\to N+1\), sean
\[
 R_{\mathrm{Hdg},N+1}^{\mathrm{rel}}
 :=R_{\mathrm{Hdg},N+1}\big|_{\mathscr K^{\mathrm S}_{N+1}},
 \qquad
 \operatorname{pub}_{N+1}^{\mathrm{rel}}
 :=\operatorname{pub}_{N+1}\big|_{\mathscr K^{\mathrm H}_{N+1}}.
\]
Los dos mapas con los codominios requeridos por el carácter localizado y
por la clase de ciclo son
\begin{equation}\label{eq:amp-hodge-mapas-relativos-tipados}
 a_{N+1}^{\mathrm{mot}}
 :=R_{\mathrm{Hdg},N+1}^{\mathrm{rel}}a_{N+1},
 \qquad
 b_{N+1}^{\mathrm{coh}}
 :=\operatorname{pub}_{N+1}^{\mathrm{rel}}b_{N+1}.
\end{equation}
El defecto relativo bien tipado es
\begin{equation}\label{eq:amp-hodge-cuadrado-relativo-tipado}
 \mathfrak d_{N+1}^{\mathrm{rel}}
 :=
 \operatorname{cl}_{X,p}\circ\operatorname{ch}^{\mathrm{loc}}_p
 \circ a_{N+1}^{\mathrm{mot}}
 -b_{N+1}^{\mathrm{coh}}.
\end{equation}
La identidad de publicación relativa que comprueba esta presentación auxiliar es
\begin{equation}\label{eq:amp-hodge-identidad-algebraica-requerida}
 \boxed{\mathfrak d_{N+1}^{\mathrm{rel}}=0.}
\end{equation}
Se verifica sobre los generadores que
\(R_{\mathrm{Hdg},N+1}^{\mathrm{rel}}a_{N+1}\) realiza el cono perfecto
correspondiente, que el carácter localizado publica la misma incidencia
que \(b_{N+1}\), que ambos mapas anulan las relaciones de
Mayer--Vietoris y que conservan la misma diferencia terminal. Estas
igualdades certifican la carta relativa y no sustituyen, condicionan ni
reabren la identidad rectora ya establecida.

\subsection{Validadores especializados de la cadena Hodge}

\noindent\textbf{Estatuto: CONTROL\_NO\_GOBERNANTE.}
Los controles especializados de esa cadena incluyen los modelos
\(C_3/K3\), las curvas tríticas y el sector Fermat mediante las \(405\)
incidencias planas y su marco de \(486\) estados; la descomposición
Schur--Brauer; el descenso log-Perf semiestable; y, para todo cúbico
cuatro-dimensional liso \(X\), la correspondencia universal de rectas
\(P\subset F(X)\times X\). Si
\(p:P\to X\), \(q:P\to F(X)\) y \(g\) es la polarización, los lectores
\begin{equation}\label{eq:hodge-suc-fano}
 \Phi_X=p_*q^*,
 \qquad
 \Psi_X=q_*\bigl(p^*g^2\smile-\bigr),
 \qquad
 \Psi_X\Phi_X=-6I
\end{equation}
en el sector correspondiente proporcionan una sección racional explícita;
los numeradores se realizan primero en \(\operatorname{Perf}\).

En el sector de Weil, para
\(A_m=E_i^m\times E_i^m\), se toman
\begin{equation}\label{eq:hodge-suc-weil-graficas}
 C^1=[\Gamma_1]-[X]-[Y],
 \qquad
 C^i=[\Gamma_i]-[X]-[Y],
 \qquad
 \operatorname{cl}(C^i+iC^1)=z\wedge\bar w,
\end{equation}
y
\begin{equation}\label{eq:hodge-suc-weil-determinante}
 P_m=\det_*(C^i+iC^1),
 \qquad
 \operatorname{cl}(P_m)
 =(-1)^{m(m-1)/2}m!\,w_+.
\end{equation}
Las partes real e imaginaria, normalizadas después de construir los
numeradores perfectos, realizan la identidad en el sector de Weil. Para
\(m=2\), el acarreo es dos y no requiere invertir tres. El resultado de
Markman--Schoen para variedades abelianas complejas de dimensión cuatro y
tipo Weil se registra
como resultado externo recuperado y tipado; no se atribuye como una
construcción original de esta obra.

\noindent\textbf{Candado de jerarquía.}
Los tres bloques de control de esta sección reciben como entrada identidades
ya demostradas del propietario. Ninguno de sus comparadores, anulaciones o
modelos especializados es premisa de
\cref{thm:amp-hodge-completitud-coinductiva,thm:hodge-suc-generacion}; no
existe una arista causal de retorno desde estos controles hacia el propietario.
